\section{Final-slice Einstein--scalar application}
\label{sec:gravity}

This appendix is not part of the paper's headline theorem.  The all-angular
de Sitter result does not require spherical symmetry.  We now restrict
to the s-wave and to a final slice on which the radial field coordinate
vanishes.  This subclass is useful because its momentum density vanishes and
its energy measure has a direct constraint interpretation.

Let $u(t,x)$ be the normalized rescaled s-wave.  The physical conformal scalar
is
\begin{equation}
 \phi(t,r)=\frac{u(t,x(r))}{\sqrt{4\pi}\,r}.
 \label{eq:physical-radial-field}
\end{equation}
On a final slice with $u=0$ and $p=\partial_tu$, the physical normal derivative
in the de Sitter geometry is
\begin{equation}
 v(r):=n^a\nabla_a\phi
 =\frac{p(x(r))}{\sqrt{4\pi}\,r\sqrt{N(r)}}.
 \label{eq:normal-derivative}
\end{equation}
For the conformally coupled stress tensor, every improvement term containing
$\phi$ or a spatial derivative of $\phi$ vanishes on this slice.  Directly,
\begin{equation}
 \rho=T_{ab}n^an^b=\frac12v^2,
 \qquad
 j_i=-T_{ab}n^ah^b{}_i=0.
 \label{eq:flux-free-stress}
\end{equation}
Since $\dd r=N\dd x$, \cref{eq:normal-derivative,eq:flux-free-stress} imply
the exact measure identity
\begin{equation}
 4\pi r^2\rho(r)\dd r=\frac12p(x)^2\dd x.
 \label{eq:mass-energy-measure}
\end{equation}

Use the same numerical function $v$ as the scalar normal-momentum datum on a
new spherical initial slice, and consider gravitational data with vanishing
extrinsic curvature,
\begin{equation}
 \dd l^2=\frac{\dd r^2}{f(r)}+r^2\dd\Omega_2^2,
 \qquad K_{ij}=0,
 \qquad
 f(r)=N(r)-\frac{2Gm(r)}{r}.
 \label{eq:constraint-metric}
\end{equation}
The momentum constraint is satisfied by \cref{eq:flux-free-stress}.  The
Hamiltonian constraint with cosmological radius $R$ is solved by
\begin{equation}
 m'(r)=4\pi r^2\rho(r),
 \qquad m(0)=0.
 \label{eq:mass-constraint}
\end{equation}
Absorbing the Bunch--Davies vacuum expectation into the background
cosmological constant, the coherent excess stress in
\cref{eq:energy-stress} supplies precisely the classical source above.
Although $K_{ij}=0$, the nonzero scalar normal momentum means that the full
Einstein--scalar data are not invariant under time reflection.
Equation~\eqref{eq:mass-energy-measure} then gives
\begin{equation}
 m(b)=\EK
 \label{eq:mass-equals-energy}
\end{equation}
for any areal radius $b$ outside the field support.

Define the dimensionless radial constraint ratio
\begin{equation}
 Q(r)=\frac{2Gm(r)}{rN(r)}.
 \label{eq:constraint-ratio}
\end{equation}
For a comparison ball $0<r\leq b<R$, set
\begin{equation}
 Q_b:=\sup_{0<r\leq b}Q(r).
\end{equation}
If $Q_b\leq\delta<1$, then throughout that ball
\begin{equation}
 f(r)\geq N(r)(1-\delta)>0,
 \qquad
 0\leq\frac{g_{rr}}{g_{rr}^{\rm dS}}-1
 \leq\frac{\delta}{1-\delta}.
 \label{eq:metric-control}
\end{equation}

\begin{corollary}[Necessary final-slice constraint cost]
\label{cor:gravity-cost}
Let the hypotheses of \cref{cor:source-bound} hold and suppose the source is
engineered to finish with spherical $q=0$ data.  Set
\begin{equation}
 b=R\tanh(L/R),
 \qquad L=R\operatorname{arctanh}(a/R)+T.
\end{equation}
Then
\begin{equation}
 Q_b\geq Q(b)
 \geq
 \frac{2G}{bN(b)R C_{\rm opt}(L/R)}
 \log\frac{1}{2\eobs}.
 \label{eq:gravity-lower-bound}
\end{equation}
Equivalently, a local budget $Q_b\leq\delta<1$ implies
\begin{equation}
 \eobs\geq\frac12\exp\left[
 -\frac{\delta\,bN(b)R}{2G}C_{\rm opt}(L/R)
 \right].
 \label{eq:gravity-error-floor}
\end{equation}
\end{corollary}

\begin{proof}
Outside the support, \cref{eq:mass-equals-energy} gives
$Q(b)=2G\EK/[bN(b)]$.  Apply the sharp energy lower bound in
\cref{cor:source-bound}; solving the same inequality in the other direction
gives \cref{eq:gravity-error-floor}.
\end{proof}

The compact mass shifts the cosmological horizon relative to the pure de
Sitter value $R$.  Accordingly, $Q_b$ is a local comparison norm on the ball
containing the final datum; no global smallness claim relative to the
unperturbed lapse is intended.

\subsection{An explicit local weak-constraint window}

For the source optical radius
$\ell=R\operatorname{arctanh}(a/R)$, take the ideal profile
\begin{equation}
 q=0,
 \qquad
 p(x)=A\sqrt{\frac{3}{\ell^3}}x\mathbf1_{(0,\ell)}(x).
 \label{eq:gravity-profile}
\end{equation}
Its cumulative mass is
\begin{equation}
 m(x)=\EK\left(\frac{x}{\ell}\right)^3.
 \label{eq:cumulative-mass}
\end{equation}
Inside the support, the function
\begin{equation}
 \frac{\operatorname{arctanh}(r/R)^3}{r[1-r^2/R^2]}
 \label{eq:q-monotone-function}
\end{equation}
is strictly increasing on $(0,a]$.  One way to see this is to factor it as
$(x/r)(x^2/N)$; both factors increase with $r$.  Hence the wall value controls
the complete profile.  Combining \cref{eq:constructive-energy} with the wall
constraint gives the sufficient estimate
\begin{equation}
 Q_a
 \leq\frac{2\pi G\Gamma_\star}
 {3\ell\,a[1-a^2/R^2]}.
 \label{eq:constructive-gravity}
\end{equation}

For illustration, set
\begin{equation}
 \eobs=10^{-6},\quad
 a/R=0.2,\quad T/R=0.1,\quad
 G/R^2=10^{-6},\quad \delta=0.25.
\end{equation}
The universal theorem requires
\begin{equation}
 C_{\rm opt}(L/R)\leq0.340766570374,
 \qquad \EK R\geq38.5083647231,
 \qquad Q(b)\geq2.86896268303\times10^{-4},
\end{equation}
whereas the constructive family has
\begin{equation}
 \EK R\leq67.7824380936,
 \qquad Q_a\leq7.06067063475\times10^{-4}.
\end{equation}
Thus the fixed numerical point has a nonempty local weak-constraint window.  By
\cref{lem:smooth-density}, smooth compact profiles preserve this strict
margin.

The construction closes only the final Hamiltonian and momentum constraints.
It does not specify the lapse, evolve the Einstein--scalar system, include the
source actuator during switching, or recompute the quantum channel in the
perturbed geometry.
